\section{Case study: a debutanizer digital twin}
\label{sec:case}

We demonstrate the certificate of Theorem~\ref{thm:composed} and the operating
rule of Corollary~\ref{cor:floor} in closed loop on a calibrated digital twin of a
debutanizer column. The three modules are the real implementations of the
constituent papers; the plant is a twin rather than an operating column. This is
a deliberate scoping choice, justified at the end of this section: the certificate
is a statement about how three coverage guarantees \emph{compose} in feedback, and
a twin that responds to the optimizer's setpoint moves is the appropriate
instrument for that question, whereas component accuracy is established on
operational data in the constituent papers.

\subsection{The plant twin}
\label{sec:twin}

The plant truth is a Fenske--Underwood--Gilliland (FUG) shortcut column model
($\alpha_{\mathrm{LK}}=5.97$, eight equilibrium stages, feed $100$~kmol\,h$^{-1}$,
light-key feed fraction $0.35$): for a decision $(R,D)$ it returns the bottoms
light-key fraction $x$, which is the unmeasured quality the soft sensor must
infer. The inferential measurement is a tray-6 temperature, modelled as cooling
monotonically with reflux as the split sharpens and landing in the
$[100.5, 111.5]\,^\circ$C window of the soft sensor's physics bridge; it carries a
$0.25\,^\circ$C sensor noise, and the bottoms truth carries a
$1.0\times10^{-3}$ process noise. The thermophysical properties (relative
volatility, heavy-key $K$-value, and liquid viscosity) are evaluated from
DIPPR-101 vapor pressures for the $n$-butane/$n$-hexane key pair.

The reflux pump degrades as a function of the optimizer's decision, closing in the
twin the endogeneity loop identified in the framework. Damage accrues at a rate
proportional to $\mathrm{load}^{p}$, where the load is the affinity brake-horsepower
ratio $\mathrm{load} = (Q/Q_{\mathrm{ref}})^{3}$ (Perry's Table~10-13) and the
exponent $p=2$ follows from rolling-element bearing fatigue (the $L_{10}$ life
varies as $\mathrm{load}^{-3}$) under a hydraulic radial load that scales as $Q^{2}$
by the affinity head law, so that damage scales as $Q^{6}$ in reflux flow. The base
rate is set to a fresh-nominal remaining life of $\approx\!1000$~h; these
coefficients are calibrated to the bearing-degradation timescale of the prognostic
module. Degradation severity drives a five-feature synthesizer whose output the
prognostic module scores.

\subsection{The three modules on the twin}
\label{sec:modules}

\emph{Module~1 (soft sensor).} The real physics-anchored soft
sensor~\citep{busico_m1} is trained on twin samples: a standardized ridge estimator
on four physics-derived features (the C4 bubble-point fraction, the relative
volatility, the stripping factor, and the delayed normalized tray temperature),
trained over $700$ operating points spanning the box and split $450$ for fitting
and $250$ for the conformal calibration residuals. On held-out twin data its adaptive conformal
interval attains empirical coverage $0.898$ against a $0.90$ target, with a mean
half-width of $7.1\times10^{-3}$, comparable to the product specification, so the
interval materially backs off the optimizer through
\eqref{eq:quality-surrogate}.

\emph{Module~2 (prognostics).} The real predictive-maintenance
stack~\citep{busico_m2} is built from its canonical configuration: a five-feature
Hotelling $T^{2}$ health index fitted on a healthy reference (alarm threshold
$T^{2}=15.1$) and a trajectory-similarity remaining-life estimator over a library of
run-to-failure arcs. Driven by the synthesizer, the health index rises
monotonically with severity and alarms at severity $\approx\!0.71$.

\emph{Module~3 (optimizer).} The real modifier-adaptation
RTO~\citep{busico_m3} optimizes economic profit over the FUG response surface
(a log-quadratic fit in $(R,D)$, coefficient of determination $0.982$), enforcing
the backed-off quality constraint \eqref{eq:quality-surrogate} and, in the
health-constrained arm, the $\psi$-budget \eqref{eq:psi-budget}. The program is
solved to local optimality at each cycle.

\subsection{Protocol}
\label{sec:protocol}

The demonstration contrasts two arms that are identical in every respect except the
$\psi$-budget. In the \emph{health-blind} arm the budget is inactive
($\varepsilon\to\infty$) and the optimizer maximizes economics subject only to the
quality surrogate. In the \emph{health-constrained} arm the budget is active with
radius $\varepsilon=0.05$. Each arm is run for $8$ random seeds of $80$ control
cycles; the coverage harness records, per cycle, whether the system-safety event
\eqref{eq:safety-event} holds against the truths, and aggregates the empirical
$S_k$-coverage with a Wilson interval. Both arms share the parameters of
Table~\ref{tab:params}. With $\alpha_1=\alpha_2=0.1$ and $\varepsilon=0.05$, the
certified floor of Corollary~\ref{cor:floor} is
$1-(\alpha_1+\alpha_2)-\varepsilon = 0.75$.

The Lipschitz constant of the soft-sensor score-quantile in $\psi$-space, estimated
a~priori from a regime sweep across the operating box, is
$L_1 \approx 0.024$. The demonstration uses a deliberately conservative common value
$L_1=L_2=0.15$ in the enforced budget \eqref{eq:psi-budget}; because this exceeds
the swept estimate, the budget is tighter than the data require and the certified
floor holds a fortiori.

\begin{table}[t]
\centering
\caption{Demonstration parameters. Both arms are identical except the
$\psi$-budget radius $\varepsilon$.}
\label{tab:params}
\begin{tabular}{lll}
\toprule
Symbol & Meaning & Value \\
\midrule
$\alpha_1$ & soft-sensor miscoverage & $0.10$ \\
$\alpha_2$ & prognostic miscoverage & $0.10$ \\
$\varepsilon$ & $\psi$-budget radius (constrained arm) & $0.05$ \\
$x^{\mathrm{spec}}$ & bottoms C4 specification & $8.0\times10^{-3}$ \\
$\rho_{\min}$ & remaining-life floor & $200$~h \\
$L_1=L_2$ & Lipschitz constants used (conservative) & $0.15$ \\
$(R,D)_{\mathrm{nom}}$ & nominal optimum & $(2.037,\,33.49)$ \\
seeds $\times$ cycles & Monte-Carlo replicates & $8\times80$ \\
\bottomrule
\end{tabular}
\end{table}

\subsection{Scope of the demonstration}
\label{sec:scope}

The closed-loop demonstration runs on a calibrated digital twin; the
\emph{accuracy} of each module is validated on operational data in the constituent
papers~\citep{busico_m1,busico_m2,busico_m3}. Two reasons make the twin the correct
instrument here. First, a closed-loop optimizer requires a plant that
\emph{responds} to its setpoint moves, because the endogeneity identified in the
framework cannot be exercised against a frozen historical dataset. Second, the
certificate concerns the \emph{composition} of three coverage guarantees under
feedback, a property of the loop architecture rather than of any module's point
accuracy; the twin isolates exactly that property while remaining fully
reproducible from open code with no proprietary data. Physical realization on a
dynamic high-fidelity simulator and on plant data is the subject of future work.
